\documentclass{beamer}
\usepackage{lmodern}

\usepackage{xeCJK}

\usepackage[safe]{silence}
\WarningFilter*{latexfont}{Font}

\usetheme{CambridgeUS} % try Madrid, Pittsburgh
\usecolortheme{beaver}
\usefonttheme[onlymath]{serif} % try "professionalfonts"

\setbeamertemplate{theorems}[normal font]
\setbeamertemplate{itemize items}[default]
\setbeamertemplate{enumerate items}[default]

\usepackage{amsmath, amsfonts, latexsym, mathtools, extarrows}
% \usepackage{kbordermatrix}
% \usepackage{ntheorem}

\usepackage{centernot}

\DeclareMathOperator*{\argmin}{arg\,min}
\DeclareMathOperator*{\argmax}{arg\,max}

% colors
\newcommand{\red}[1]{\textcolor{red}{#1}}
\newcommand{\green}[1]{\textcolor{green}{#1}}
\newcommand{\blue}[1]{\textcolor{blue}{#1}}
\newcommand{\purple}[1]{\textcolor{purple}{#1}}

% color box
\newcommand{\rbox}[1]{\red{\boxed{#1}}}
\newcommand{\gbox}[1]{\green{\boxed{#1}}}
\newcommand{\bbox}[1]{\blue{\boxed{#1}}}
\newcommand{\pbox}[1]{\purple{\boxed{#1}}}

\usepackage{tikz}
% see http://tex.stackexchange.com/a/7045/23098
\newcommand*\circled[1]{\tikz[baseline=(char.base)]{
            \node[shape=circle,draw,inner sep=2pt] (char) {#1};}}
\usepackage{tikz-qtree}
\usetikzlibrary{backgrounds, fit}

\usepackage{pifont}
\usepackage{wasysym}

\usepackage[normalem]{ulem}
\newcommand{\middlewave}[1]{\raisebox{0.5em}{\uwave{\hspace{#1}}}}

\usepackage{graphicx, subcaption}

\usepackage{algorithm}
\usepackage[noend]{algpseudocode}

\newcommand{\pno}[1]{\textcolor{blue}{\scriptsize [Problem: #1]}}
\newcommand{\set}[1]{\{#1\}}
\DeclareMathOperator{\lcm}{lcm}

\newcommand{\cmark}{\green{\ding{51}}}
\newcommand{\xmark}{\red{\ding{55}}}
%%%%%%%%%%%%%%%%%%%%%%%%%%%%%%%%%%%%%%%%%%%%%%%%%%%%%%%%%%%%%%
% for fig without caption: #1: width/size; #2: fig file
\newcommand{\fignocaption}[2]{
  \begin{figure}[htp]
    \centering
      \includegraphics[#1]{#2}
  \end{figure}
}

% for fig with caption: #1: width/size; #2: fig file; #3: fig caption
\newcommand{\fig}[3]{
  \begin{figure}[htp]
    \centering
      \includegraphics[#1]{#2}
      \caption[labelInTOC]{#3}
  \end{figure}
}

\newcommand{\p}{\text{P}}
\newcommand{\np}{\text{NP}}
\newcommand{\conp}{\text{coNP}}
\newcommand{\pr}{\le_{p}} % polynomial reduction
\newcommand{\nph}{\text{NP-hard}}
\newcommand{\npc}{\text{NP-complete}}
\newcommand{\npi}{\text{NP-intermediate}}
\newcommand{\conpc}{\text{coNP-complete}}
\newcommand{\expcls}{\text{EXP}}
\newcommand{\rcls}{\text{R}}
\newcommand{\dtime}{\text{DTIME}}
\newcommand{\dspace}{\text{SPACE}}
\newcommand{\pspace}{\text{PSPACE}}
\newcommand{\pspacec}{\text{PSPACE-complete}}
\newcommand{\ph}{\text{PH}}
\newcommand{\hp}{\text{HAM-PATH}}
\newcommand{\hc}{\text{HAM-CYCLE}}
\newcommand{\fp}{\text{FP}}
\newcommand{\poly}{\text{Poly}}
\newcommand{\po}{\text{PO}}
\newcommand{\fptas}{\text{FPTAS}}
\newcommand{\ptas}{\text{PTAS}}
\newcommand{\apx}{\text{APX}}
\newcommand{\logapx}{\text{Log-APX}}
\newcommand{\polyapx}{\text{Poly-APX}}
\newcommand{\expapx}{\text{Exp-APX}}
\newcommand{\npo}{\text{NPO}}
\newcommand{\ms}{\text{MS}}

\newcommand{\titletext}{Approximation Algorithms}

\newcommand{\tx}{%
  \section*{}
  \begin{frame}[noframenumbering]
	\fignocaption{width = 0.50\textwidth}{figs/thankyou.jpg}
  \end{frame}
}
%%%%%%%%%%%%%%%%%%%%
\title[\titletext]{\titletext}
% \subtitle{}

\author[Hengfeng Wei]{Hengfeng Wei}
\titlegraphic{\includegraphics[height = 1.5cm]{figs/qrcode-approx-alg-20170518.png}}
\institute{hfwei@nju.edu.cn}
\date{May 14 $\sim$ May 18, 2017}

\AtBeginSection[]{
  \begin{frame}[noframenumbering, plain]
    \frametitle{\titletext}
    \tableofcontents[currentsection, sectionstyle=show/shaded, subsectionstyle=show/show/hide]
  \end{frame}
}
%%%%%%%%%%
\begin{document}
\maketitle

\section{Approximation Classes}

\documentclass{beamer}
\usepackage{lmodern}

\usepackage{xeCJK}

\usepackage[safe]{silence}
\WarningFilter*{latexfont}{Font}

\usetheme{CambridgeUS} % try Madrid, Pittsburgh
\usecolortheme{beaver}
\usefonttheme[onlymath]{serif} % try "professionalfonts"

\setbeamertemplate{theorems}[normal font]
\setbeamertemplate{itemize items}[default]
\setbeamertemplate{enumerate items}[default]

\usepackage{amsmath, amsfonts, latexsym, mathtools, extarrows}
% \usepackage{kbordermatrix}
% \usepackage{ntheorem}

\usepackage{centernot}

\DeclareMathOperator*{\argmin}{arg\,min}
\DeclareMathOperator*{\argmax}{arg\,max}

% colors
\newcommand{\red}[1]{\textcolor{red}{#1}}
\newcommand{\green}[1]{\textcolor{green}{#1}}
\newcommand{\blue}[1]{\textcolor{blue}{#1}}
\newcommand{\purple}[1]{\textcolor{purple}{#1}}

% color box
\newcommand{\rbox}[1]{\red{\boxed{#1}}}
\newcommand{\gbox}[1]{\green{\boxed{#1}}}
\newcommand{\bbox}[1]{\blue{\boxed{#1}}}
\newcommand{\pbox}[1]{\purple{\boxed{#1}}}

\usepackage{tikz}
% see http://tex.stackexchange.com/a/7045/23098
\newcommand*\circled[1]{\tikz[baseline=(char.base)]{
            \node[shape=circle,draw,inner sep=2pt] (char) {#1};}}
\usepackage{tikz-qtree}
\usetikzlibrary{backgrounds, fit}

\usepackage{pifont}
\usepackage{wasysym}

\usepackage[normalem]{ulem}
\newcommand{\middlewave}[1]{\raisebox{0.5em}{\uwave{\hspace{#1}}}}

\usepackage{graphicx, subcaption}

\usepackage{algorithm}
\usepackage[noend]{algpseudocode}

\newcommand{\pno}[1]{\textcolor{blue}{\scriptsize [Problem: #1]}}
\newcommand{\set}[1]{\{#1\}}
\DeclareMathOperator{\lcm}{lcm}

\newcommand{\cmark}{\green{\ding{51}}}
\newcommand{\xmark}{\red{\ding{55}}}
%%%%%%%%%%%%%%%%%%%%%%%%%%%%%%%%%%%%%%%%%%%%%%%%%%%%%%%%%%%%%%
% for fig without caption: #1: width/size; #2: fig file
\newcommand{\fignocaption}[2]{
  \begin{figure}[htp]
    \centering
      \includegraphics[#1]{#2}
  \end{figure}
}

% for fig with caption: #1: width/size; #2: fig file; #3: fig caption
\newcommand{\fig}[3]{
  \begin{figure}[htp]
    \centering
      \includegraphics[#1]{#2}
      \caption[labelInTOC]{#3}
  \end{figure}
}

\newcommand{\p}{\text{P}}
\newcommand{\np}{\text{NP}}
\newcommand{\conp}{\text{coNP}}
\newcommand{\pr}{\le_{p}} % polynomial reduction
\newcommand{\nph}{\text{NP-hard}}
\newcommand{\npc}{\text{NP-complete}}
\newcommand{\npi}{\text{NP-intermediate}}
\newcommand{\conpc}{\text{coNP-complete}}
\newcommand{\expcls}{\text{EXP}}
\newcommand{\rcls}{\text{R}}
\newcommand{\dtime}{\text{DTIME}}
\newcommand{\dspace}{\text{SPACE}}
\newcommand{\pspace}{\text{PSPACE}}
\newcommand{\pspacec}{\text{PSPACE-complete}}
\newcommand{\ph}{\text{PH}}
\newcommand{\hp}{\text{HAM-PATH}}
\newcommand{\hc}{\text{HAM-CYCLE}}
\newcommand{\partition}{\text{3-PARTITION}}
\newcommand{\tetris}{\text{TETRIS}}

\newcommand{\titletext}{Approximation Algorithms}

\newcommand{\tx}{%
  \section*{}
  \begin{frame}[noframenumbering]
	\fignocaption{width = 0.50\textwidth}{figs/thankyou.jpg}
  \end{frame}
}
%%%%%%%%%%%%%%%%%%%%
\title[\titletext]{\titletext}
% \subtitle{}

\author[Hengfeng Wei]{Hengfeng Wei}
% \titlegraphic{\includegraphics[height = 1.5cm]{figs/qrcode-approx-alg-20170518.png}}
\institute{hfwei@nju.edu.cn}
\date{May 14 $\sim$ May 18, 2017}

\AtBeginSection[]{
  \begin{frame}[noframenumbering, plain]
    \frametitle{\titletext}
    \tableofcontents[currentsection, sectionstyle=show/shaded, subsectionstyle=show/show/hide]
  \end{frame}
}
%%%%%%%%%%
\begin{document}
\maketitle

\section{Approximation Classes}

\documentclass{beamer}
\usepackage{lmodern}

\usepackage{xeCJK}

\usepackage[safe]{silence}
\WarningFilter*{latexfont}{Font}

\usetheme{CambridgeUS} % try Madrid, Pittsburgh
\usecolortheme{beaver}
\usefonttheme[onlymath]{serif} % try "professionalfonts"

\setbeamertemplate{theorems}[normal font]
\setbeamertemplate{itemize items}[default]
\setbeamertemplate{enumerate items}[default]

\usepackage{amsmath, amsfonts, latexsym, mathtools, extarrows}
% \usepackage{kbordermatrix}
% \usepackage{ntheorem}

\usepackage{centernot}

\DeclareMathOperator*{\argmin}{arg\,min}
\DeclareMathOperator*{\argmax}{arg\,max}

% colors
\newcommand{\red}[1]{\textcolor{red}{#1}}
\newcommand{\green}[1]{\textcolor{green}{#1}}
\newcommand{\blue}[1]{\textcolor{blue}{#1}}
\newcommand{\purple}[1]{\textcolor{purple}{#1}}

% color box
\newcommand{\rbox}[1]{\red{\boxed{#1}}}
\newcommand{\gbox}[1]{\green{\boxed{#1}}}
\newcommand{\bbox}[1]{\blue{\boxed{#1}}}
\newcommand{\pbox}[1]{\purple{\boxed{#1}}}

\usepackage{tikz}
% see http://tex.stackexchange.com/a/7045/23098
\newcommand*\circled[1]{\tikz[baseline=(char.base)]{
            \node[shape=circle,draw,inner sep=2pt] (char) {#1};}}
\usepackage{tikz-qtree}
\usetikzlibrary{backgrounds, fit}

\usepackage{pifont}
\usepackage{wasysym}

\usepackage[normalem]{ulem}
\newcommand{\middlewave}[1]{\raisebox{0.5em}{\uwave{\hspace{#1}}}}

\usepackage{graphicx, subcaption}

\usepackage{algorithm}
\usepackage[noend]{algpseudocode}

\newcommand{\pno}[1]{\textcolor{blue}{\scriptsize [Problem: #1]}}
\newcommand{\set}[1]{\{#1\}}
\DeclareMathOperator{\lcm}{lcm}

\newcommand{\cmark}{\green{\ding{51}}}
\newcommand{\xmark}{\red{\ding{55}}}
%%%%%%%%%%%%%%%%%%%%%%%%%%%%%%%%%%%%%%%%%%%%%%%%%%%%%%%%%%%%%%
% for fig without caption: #1: width/size; #2: fig file
\newcommand{\fignocaption}[2]{
  \begin{figure}[htp]
    \centering
      \includegraphics[#1]{#2}
  \end{figure}
}

% for fig with caption: #1: width/size; #2: fig file; #3: fig caption
\newcommand{\fig}[3]{
  \begin{figure}[htp]
    \centering
      \includegraphics[#1]{#2}
      \caption[labelInTOC]{#3}
  \end{figure}
}

\newcommand{\p}{\text{P}}
\newcommand{\np}{\text{NP}}
\newcommand{\conp}{\text{coNP}}
\newcommand{\pr}{\le_{p}} % polynomial reduction
\newcommand{\nph}{\text{NP-hard}}
\newcommand{\npc}{\text{NP-complete}}
\newcommand{\npi}{\text{NP-intermediate}}
\newcommand{\conpc}{\text{coNP-complete}}
\newcommand{\expcls}{\text{EXP}}
\newcommand{\rcls}{\text{R}}
\newcommand{\dtime}{\text{DTIME}}
\newcommand{\dspace}{\text{SPACE}}
\newcommand{\pspace}{\text{PSPACE}}
\newcommand{\pspacec}{\text{PSPACE-complete}}
\newcommand{\ph}{\text{PH}}
\newcommand{\hp}{\text{HAM-PATH}}
\newcommand{\hc}{\text{HAM-CYCLE}}
\newcommand{\partition}{\text{3-PARTITION}}
\newcommand{\tetris}{\text{TETRIS}}

\newcommand{\titletext}{Approximation Algorithms}

\newcommand{\tx}{%
  \section*{}
  \begin{frame}[noframenumbering]
	\fignocaption{width = 0.50\textwidth}{figs/thankyou.jpg}
  \end{frame}
}
%%%%%%%%%%%%%%%%%%%%
\title[\titletext]{\titletext}
% \subtitle{}

\author[Hengfeng Wei]{Hengfeng Wei}
% \titlegraphic{\includegraphics[height = 1.5cm]{figs/qrcode-approx-alg-20170518.png}}
\institute{hfwei@nju.edu.cn}
\date{May 14 $\sim$ May 18, 2017}

\AtBeginSection[]{
  \begin{frame}[noframenumbering, plain]
    \frametitle{\titletext}
    \tableofcontents[currentsection, sectionstyle=show/shaded, subsectionstyle=show/show/hide]
  \end{frame}
}
%%%%%%%%%%
\begin{document}
\maketitle

\section{Approximation Classes}

\documentclass{beamer}
\usepackage{lmodern}

\usepackage{xeCJK}

\usepackage[safe]{silence}
\WarningFilter*{latexfont}{Font}

\usetheme{CambridgeUS} % try Madrid, Pittsburgh
\usecolortheme{beaver}
\usefonttheme[onlymath]{serif} % try "professionalfonts"

\setbeamertemplate{theorems}[normal font]
\setbeamertemplate{itemize items}[default]
\setbeamertemplate{enumerate items}[default]

\usepackage{amsmath, amsfonts, latexsym, mathtools, extarrows}
% \usepackage{kbordermatrix}
% \usepackage{ntheorem}

\usepackage{centernot}

\DeclareMathOperator*{\argmin}{arg\,min}
\DeclareMathOperator*{\argmax}{arg\,max}

% colors
\newcommand{\red}[1]{\textcolor{red}{#1}}
\newcommand{\green}[1]{\textcolor{green}{#1}}
\newcommand{\blue}[1]{\textcolor{blue}{#1}}
\newcommand{\purple}[1]{\textcolor{purple}{#1}}

% color box
\newcommand{\rbox}[1]{\red{\boxed{#1}}}
\newcommand{\gbox}[1]{\green{\boxed{#1}}}
\newcommand{\bbox}[1]{\blue{\boxed{#1}}}
\newcommand{\pbox}[1]{\purple{\boxed{#1}}}

\usepackage{tikz}
% see http://tex.stackexchange.com/a/7045/23098
\newcommand*\circled[1]{\tikz[baseline=(char.base)]{
            \node[shape=circle,draw,inner sep=2pt] (char) {#1};}}
\usepackage{tikz-qtree}
\usetikzlibrary{backgrounds, fit}

\usepackage{pifont}
\usepackage{wasysym}

\usepackage[normalem]{ulem}
\newcommand{\middlewave}[1]{\raisebox{0.5em}{\uwave{\hspace{#1}}}}

\usepackage{graphicx, subcaption}

\usepackage{algorithm}
\usepackage[noend]{algpseudocode}

\newcommand{\pno}[1]{\textcolor{blue}{\scriptsize [Problem: #1]}}
\newcommand{\set}[1]{\{#1\}}
\DeclareMathOperator{\lcm}{lcm}

\newcommand{\cmark}{\green{\ding{51}}}
\newcommand{\xmark}{\red{\ding{55}}}
%%%%%%%%%%%%%%%%%%%%%%%%%%%%%%%%%%%%%%%%%%%%%%%%%%%%%%%%%%%%%%
% for fig without caption: #1: width/size; #2: fig file
\newcommand{\fignocaption}[2]{
  \begin{figure}[htp]
    \centering
      \includegraphics[#1]{#2}
  \end{figure}
}

% for fig with caption: #1: width/size; #2: fig file; #3: fig caption
\newcommand{\fig}[3]{
  \begin{figure}[htp]
    \centering
      \includegraphics[#1]{#2}
      \caption[labelInTOC]{#3}
  \end{figure}
}

\newcommand{\p}{\text{P}}
\newcommand{\np}{\text{NP}}
\newcommand{\conp}{\text{coNP}}
\newcommand{\pr}{\le_{p}} % polynomial reduction
\newcommand{\nph}{\text{NP-hard}}
\newcommand{\npc}{\text{NP-complete}}
\newcommand{\npi}{\text{NP-intermediate}}
\newcommand{\conpc}{\text{coNP-complete}}
\newcommand{\expcls}{\text{EXP}}
\newcommand{\rcls}{\text{R}}
\newcommand{\dtime}{\text{DTIME}}
\newcommand{\dspace}{\text{SPACE}}
\newcommand{\pspace}{\text{PSPACE}}
\newcommand{\pspacec}{\text{PSPACE-complete}}
\newcommand{\ph}{\text{PH}}
\newcommand{\hp}{\text{HAM-PATH}}
\newcommand{\hc}{\text{HAM-CYCLE}}
\newcommand{\partition}{\text{3-PARTITION}}
\newcommand{\tetris}{\text{TETRIS}}

\newcommand{\titletext}{Approximation Algorithms}

\newcommand{\tx}{%
  \section*{}
  \begin{frame}[noframenumbering]
	\fignocaption{width = 0.50\textwidth}{figs/thankyou.jpg}
  \end{frame}
}
%%%%%%%%%%%%%%%%%%%%
\title[\titletext]{\titletext}
% \subtitle{}

\author[Hengfeng Wei]{Hengfeng Wei}
% \titlegraphic{\includegraphics[height = 1.5cm]{figs/qrcode-approx-alg-20170518.png}}
\institute{hfwei@nju.edu.cn}
\date{May 14 $\sim$ May 18, 2017}

\AtBeginSection[]{
  \begin{frame}[noframenumbering, plain]
    \frametitle{\titletext}
    \tableofcontents[currentsection, sectionstyle=show/shaded, subsectionstyle=show/show/hide]
  \end{frame}
}
%%%%%%%%%%
\begin{document}
\maketitle

\input{sections/approx-class}
\input{sections/approx-alg}
\tx{}

\end{document}
%%%%%%%%%%

\tx{}

\end{document}
%%%%%%%%%%

\tx{}

\end{document}
%%%%%%%%%%

\tx{}

\end{document}
%%%%%%%%%%
